% ============================================================
% BACKGROUND — Full version for NeurIPS 2026
% ============================================================

\section{Background}
\label{sec:background}

\subsection{State Space Models and Mamba}

A continuous-time linear state space model (SSM) maps an input signal
$x(t) \in \R$ to an output $y(t) \in \R$ through a latent state
$h(t) \in \R^N$:
\begin{align}
    h'(t) &= \bA h(t) + \bB x(t), \label{eq:ssm_continuous_h} \\
    y(t) &= \bC h(t) + D x(t), \label{eq:ssm_continuous_y}
\end{align}
where $\bA \in \R^{N \times N}$ is the state transition matrix,
$\bB \in \R^{N \times 1}$ the input projection, $\bC \in \R^{1 \times N}$
the output projection, and $D \in \R$ a skip connection. After
zero-order hold (ZOH) discretization with step size $\Delta$, the
recurrence becomes:
\begin{align}
    h_t &= \bar{\bA} h_{t-1} + \bar{\bB} x_t, \label{eq:ssm_discrete_h} \\
    y_t &= \bC h_t, \label{eq:ssm_discrete_y}
\end{align}
where $\bar{\bA} = \exp(\Delta \bA)$ and $\bar{\bB} = (\Delta \bA)^{-1}
(\exp(\Delta \bA) - I) \cdot \Delta \bB$. This recurrence can be
unrolled as a convolution for efficient parallel training, while
retaining the recurrent form for fast autoregressive
inference~\citep{gu2021hippo}.

\paragraph{Selective SSMs (Mamba).}
Classical SSMs use fixed $\bA$, $\bB$, $\bC$ matrices across all
inputs. The Mamba architecture~\citep{gu2023mamba} introduces
\emph{selectivity}: at each timestep $t$, the parameters $\bB_t$,
$\bC_t$, and $\Delta_t$ are computed as learned linear projections of
the input $x_t$. This allows the model to dynamically gate what
information enters and exits the hidden state---analogous to how LSTM
gates control memory. Selectivity is what enables Mamba to model
context-dependent dynamics, making it well-suited for non-stationary
time series.

\paragraph{Multivariate formulation.}
For a $K$-dimensional time series $\bX_t \in \R^K$, we use a Mamba
model with an input projection $\bW_{\text{in}} \in \R^{d \times K}$
that maps the $K$ input channels to a $d$-dimensional model space, and
an output projection $\bW_{\text{out}} \in \R^{K \times d}$ that maps
back. The full forward computation is:
\begin{equation}
    \hat{\bX}_{t+1} = \bW_{\text{out}} \cdot \text{SSM}_L
    \big( \cdots \text{SSM}_1(\bW_{\text{in}} \cdot \bX_t) \cdots \big),
    \label{eq:full_forward}
\end{equation}
where $\text{SSM}_\ell$ denotes the $\ell$-th selective SSM layer with
residual connections and layer normalization. The key structural observation is that inter-variable information
flows through $\bW_{\text{in}}$ and $\bW_{\text{out}}$: variable $j$'s
influence on variable $i$'s prediction is primarily mediated by column
$j$ of $\bW_{\text{in}}$ and row $i$ of $\bW_{\text{out}}$. While
residual connections and layer normalization create additional pathways,
the dominant signal flow passes through these projections---a property
we verify empirically by comparing the weight-projection score against
the full Jacobian score (Section~\ref{sec:ablation}).

\subsection{Granger Causality}

\begin{definition}[Granger Causality~\citep{granger1969investigating}]
\label{def:granger}
Variable $X_j$ \emph{Granger-causes} variable $X_i$ if the past values
of $X_j$ contain information that helps predict $X_i$ beyond what is
contained in the past of all other variables, including $X_i$ itself.
\end{definition}

For a vector autoregressive process of order $p$, VAR($p$):
\begin{equation}
    \bX_t = \sum_{\tau=1}^{p} \Phi_\tau \bX_{t-\tau} + \epsilon_t,
    \quad \epsilon_t \sim \mathcal{N}(0, \Sigma),
    \label{eq:var_p}
\end{equation}
Granger causality from $X_j$ to $X_i$ is equivalent to the existence of
some lag $\tau \in \{1, \ldots, p\}$ such that $[\Phi_\tau]_{i,j} \neq 0$.
This equivalence makes VAR processes the canonical testbed for causal
discovery methods.

\subsection{Neural Granger Causality}

Tank et al.~\citep{tank2021neural} proposed learning causal graphs by
training one LSTM (or MLP) per target variable $X_i$, with group-lasso
regularization ($\ell_{2,1}$-norm) on the input weights:
\begin{equation}
    \mathcal{L}_{\text{cLSTM}} = \underbrace{\frac{1}{T} \sum_t
    \|\hat{X}_i(t) - X_i(t)\|^2}_{\text{prediction loss}} +
    \lambda \underbrace{\sum_{j=1}^{K}
    \|\bW^{(i)}_{\text{in}, :, j}\|_2}_{\text{group-lasso penalty}},
    \label{eq:clstm_loss}
\end{equation}
where $\bW^{(i)}_{\text{in}, :, j}$ is the group of input weights
connecting variable $j$ to the network predicting variable $i$. When
this group is driven to zero by regularization, $X_j$ is deemed
non-causal for $X_i$. This approach requires: (1)~training $K$ separate
networks, (2)~a causal-specific loss with tuned $\lambda$, and (3)~a
group-lasso solver. Our method eliminates all three requirements.
